% AI systems are going to be making decisions that affect the lives of the people. The way these decisions are made will be decided by the developres. It is not fair for us to let the developers of these systems decide which values are the correct ones. Therefore we need to ensure that the values of the populations that these systems are deployed in are reflected in the behaviour of these systems. 

The creation of these frontier AI systems is not a neutral process. The creators of any system imbues it with their own assumption, values and biases. Given the power of these fronter AI systems, this imbuing can have a profound impact on the systems' behaviour. The process of imbuing one set of values into a system doesn't need to be an explicit process and instead results from subtle design choices, principles and assumptions made during the creationary process. Regardless of how the creators imbue their values into the system, the result is a system with behaviours that are aligned—to some extent—to their creators.

The world is very diverse in terms of the values. Out of all these diverse set of values it is impossible to create a system that perfectly align with all of them. This is because many values are explicitly in tension with each other. Therefore one must accept that any system will be aligned to some set of values and not others. 

% All values are important
If each system is only going to aligned to a certain set of values what should those values be? It doesn't seem very just or fair to prescribe a single set of values as the \scare{correct} ones. We can look at this unfairness through the lens of values being a set of trade-offs and so the performance of different value systems based on which one optimises the survival and flourishing of the widest range of things. The challenge is evaluating this performance is hard. Rather than solve that problem we can instead acknowledge that each distinct value system is a potential solution to the question of ``which values are the best values''. The diversity should be maintained due to the ignorance of which value system is the best. 

Globally, the development and deployment of these systems require immense resources, which means that only a small number of well-funded organisations are capable of training frontier models \cite{cottierRisingCostsTraining2024}. A future in which the economic utility of these systems increase yet the number of groups that can develop them remains small, creates a situation in which much of the world has to be the recipients of systems that they themselves did not build. Given the diverse set of values that these systems will have creates much misalignment of values between the systems that are deployed and the populations that are using them.

This is a problem for two reasons. Firstly because it will slowly corrode the values of the users which is as per earlier a risk fo loosing a potentially superior value system. Secondly the misalignment of values and thus the inability for populations to shape the future based on their values is a loss of agency and autonomy. The idea of having agency and autonomy is a core value of many value systems and so is widely considered a negative outcome.

In New Zealand we are one of the smaller countries that are currently unable to develop frontier AI systems. We are going to be a country that inherits the values of the systems that are deployed in our country. Given we are already one of the most sceptical and concerned about AI in the developed world \cite{TrustAttitudesUse,GlobalAttitudesAI2026} this is a problem. New Zealand also has comparatively limited AI-specific governance: it relies on a light-touch, principles-based approach and existing laws rather than a standalone AI Act \cite{erdelyiAIRegulationStill2024,NewZealandCreated}. Therefore for New Zealand to ensure that the AI systems deployed in our country are aligned with the values of our population we need to develop low-resource methods to align these systems with the values of our population. This project is a step towards that goal.